\section{Introduction}
\label{sec:introduction}
Earthquake-triggered landslides can affect inaccessible terrain before field surveys or cloud-free optical images become available. Synthetic aperture radar (SAR) supports rapid response through cloud cover and independent of daylight \cite{Mondini2021,Handwerger2022,Amatya2023}. Convolutional and attention-based networks have advanced SAR landslide mapping \cite{Nava2022DL,Nava2022Attention,Boehm2022,Chen2024MFEWF}. Nevertheless, foreshortening, layover, and shadow alter the terrain evidence visible from each radar look direction \cite{Shi2024}; backscatter and speckle also complicate local landslide signatures \cite{Santangelo2022}. Ascending and descending observations therefore provide complementary, rather than interchangeable, views. Requiring both post-event passes can delay the first response map. We instead use paired orbits during training and only the available target orbit during deployment (Fig.~\ref{fig:scenario}).

Learning with privileged information \cite{LopezPaz2016} and knowledge distillation (KD) \cite{Hinton2015} provide a basis for transferring training-only observations. Remote-sensing approaches transfer grouped features \cite{Li2022}, multimodal online knowledge \cite{Liu2024}, or frequency and cross-modal affinity information \cite{Sun2027}. Other methods preserve modality-specific information through contrastive distillation \cite{Sikdar2025} or compensate missing features using disentanglement and classwise attention \cite{Kieu2026DIS2}. These approaches motivate a different question for view-dependent SAR: how can an unavailable complementary observation teach a model to exploit evidence that remains observable in the target view?

Deformable convolution and attention separate sampling locations and aggregation weights from the sampled values \cite{Dai2017,Zhu2021}. Adaptive aggregation has been applied to remote-sensing segmentation and road extraction \cite{Zuo2022,Dai2023}, while deformable cross-attention addresses weak alignment across imaging modalities \cite{Hu2025}. We use complementary context to guide a target-side search policy; the searched values always come from the target orbit. This prevents a dual-orbit value-fusion path that the deployed Student could not reproduce.

The key distinction is the information being transferred. A fused Teacher representation may contain measurements visible only from the complementary orbit. Matching that representation can therefore ask the Student to reproduce evidence absent from its input. Our Teacher instead determines a policy for reading target features. Applying Teacher and Student policies to the same target memory makes their discrepancy concern the retrieval operation rather than a difference between modality feature spaces. Complementary supervision can change how available context is used, but does not restore a radar measurement that was never observed.

Our contributions are a paired-orbit training framework with target-only deployment; a counter-guided search Teacher that separates policy generation from target evidence; and policy and shared-memory evidence distillation for a target-only Student. Experiments compare several distillation methods on valid and full image domains and examine complementary guidance, sampling, and loss objectives. Geometric masks are used after prediction to distinguish improvements on visible terrain from behavior under distortion; training never discards the latter pixels.
